\documentclass{scrartcl}

\usepackage{blindtext}% for dummy text

\usepackage{clao}
\setmainlanguage{english}
\addbibresource{clao-sample.bib}
\setmainfont[Mapping=tex-text, Ligatures=Common, Numbers=OldStyle, SmallCapsFeatures={Letters=UppercaseSmallCaps, WordSpace=2, LetterSpace=8}]{Linux Libertine O}% The Brill font should be used instead, but it is not available on Overleaf

%%%%% Metadata
\title{Title of the article}
\subtitle{Subtitle of the article}
%%% The short title an author are printed in the page headers
\shorttitle{Short title}
\shortauthor{Short Author}
\date{\today}
%%% Doi is optional
\doi{000111222333}
%%% Authors should be input with the following commands
\claoauthor{First Author}{Institution}{address@mail.mail}
\claoauthor{Second Author}{Institution}{address@mail.mail}
%%%%%

\begin{document}
\maketitle

\begin{abstract}
The \texttt{clao.sty} package approximates the style sheet of the journal \textit{Cahiers de Linguistique Asie Orientale} (\url{https://brill.com/clao}). Submissions to the journal are not required to use this package. The extension is however useful for getting an idea of the approximate number of pages of a manuscript and for anticipating problems with large tables and figures.

The extension can only be used with the \texttt{scrartcl} class. The following packages are automatically loaded: \texttt{calc, geometry, scrlayer-scrpage, hyperref, enumitem, datatool, gb4e, xpatch, csquotes, polyglossia, biblatex}.

The journal uses the Brill font, which can be freely downloaded from \url{https://brill.com/page/290?lang=en}, but since it is not available directly on Overleaf, this document uses the Linux Libertine O font instead.\end{abstract}

\begin{keywords}
keyword -- keyword -- keyword -- keyword -- keyword
\end{keywords}

\begin{abstractfr}
\blindtext
\end{abstractfr}

\begin{keywords}
mot clé -- mot clé -- mot clé -- mot clé -- mot clé
\end{keywords}

\begin{exe}
\ex \gll Ceci est un test\\
This is a test\\
\glt \enquote*{This is a test}
\ex \gll Ceci est un autre test\\
This is a other test\\
\glt \enquote*{This is another test}
\end{exe}

\begin{exe}
\ex
\begin{xlist}
\ex \gll Ceci est un test\\
This is a test\\
\glt \enquote*{This is a test}
\ex \gll Ceci est un autre test\\
This is a other test\\
\glt \enquote*{This is another test}
\end{xlist}
\end{exe}

\blinddocument

\nocite{*}
\printbibliography
\end{document}